\begin{frame}{지니불순도의 두 가지 해석}
  \begin{columns}[T]
    \begin{column}{0.55\textwidth}
      \kb{(1) ``잘못 고를 확률'' 해석}
      \begin{itemize}
        \item 노드에서 항목을 \textbf{무작위 2개} 뽑아 \textbf{다른 클래스}일
          확률이 정확히 $G$
        \item 이진: $1 - p^2 - (1-p)^2 = 2p(1-p)$
        \item 순수($p=1$) $\to$ ``다를 확률 0'', 반반($p=0.5$) $\to$ 최대
          0.5
        \item 즉 $G$ = ``이 노드의 예측이 \textbf{틀릴 확률}''
      \end{itemize}
      \vskip0.5em
      \kb{(2) ``스무고개 정보'' 해석}
      \begin{itemize}
        \item 질문이 ``클래스1인가?''일 때, 답이 \textbf{반반으로 갈릴수록}
          정보가 가장 크다
        \item $G = 2p(1-p)$는 $p=0.5$에서 최대인 포물선
      \end{itemize}
    \end{column}
    \begin{column}{0.42\textwidth}
      \centering
      \includegraphics[width=\linewidth]{ch07_1_gini_entropy_curves.png}
      \vskip0.3em
      {\footnotesize 둘 다 $p=0.5$에서 최대, 순수에서 0.
      지니는 매끄러운 포물선, 엔트로피는 반반 근처가 더 완만.}
    \end{column}
  \end{columns}
\end{frame}
